\chapter{Considerações Finais}

Este trabalho teve como objetivo investigar a possibilidade de integrar Modelos de Linguagem de Grande Escala, ferramentas de Análise Estática de Segurança e estratégias de contextualização para apoiar a correção automatizada de vulnerabilidades em código-fonte. A motivação central da pesquisa parte do fato de que LLMs podem auxiliar na geração de correções, mas ainda não oferecem garantias suficientes quando utilizadas de forma isolada. Da mesma forma, ferramentas SAST são úteis para identificar possíveis falhas, mas seus alertas precisam ser interpretados, validados e tratados de maneira adequada.

Ao longo do TCC 1, foi construída a fundamentação teórica necessária para sustentar a proposta, abordando LLMs, vulnerabilidades de software, OWASP Top 10, engenharia de prompt, RAG, \textit{Chain-of-Thought} e ferramentas SAST. Essa base permitiu delimitar o problema de pesquisa e orientar a construção de um pipeline experimental voltado à correção de \textit{issues} de segurança identificadas por análise estática.

A metodologia definida para o TCC 2 passou a se concentrar na aplicação do pipeline em um projeto \textit{open source}. Essa decisão torna o estudo mais próximo de um cenário prático de desenvolvimento, pois as vulnerabilidades analisadas não serão inseridas manualmente, mas identificadas a partir dos alertas emitidos pela ferramenta SAST no próprio projeto selecionado. Dessa forma, o trabalho deixa de depender de um conjunto previamente controlado de falhas e passa a avaliar o comportamento da abordagem em uma base de código real, com estrutura, dependências, testes e limitações próprias.

A prova de conceito desenvolvida nesta etapa teve papel importante na validação inicial da proposta. Ela demonstrou a viabilidade técnica de automatizar a execução de uma ferramenta SAST, coletar diagnósticos, montar diferentes estratégias de \textit{prompt}, enviar requisições à LLM, aplicar correções geradas, executar testes automatizados e consolidar métricas de avaliação. Embora tenha sido realizada com casos reduzidos e controlados, a PoC permitiu observar aspectos relevantes para a continuidade do trabalho, como a importância da rastreabilidade dos artefatos, da padronização das respostas da LLM e da combinação entre reanálise estática e testes funcionais.

Assim, o TCC 1 estabelece as bases para a continuidade da pesquisa. A principal contribuição desta etapa está na definição de um fluxo estruturado para avaliar correções geradas por LLMs de forma mais rigorosa, evitando falsas correções. A continuidade do trabalho permitirá verificar até que ponto a integração entre SAST, LLMs e estratégias de contextualização pode contribuir para apoiar a correção de vulnerabilidades de maneira segura, rastreável e funcional.